\section{Asymptotic links with Wasserstein ABC}
\label{sec:asymptotics}

In this section we clarify the relation between \textsc{permABC} and Wasserstein ABC (WABC),
and we discuss two asymptotic regimes that seem especially relevant for our framework:
(i) the number of observed compartments $K \to \infty$, and (ii) the over-sampling regime
$M \to \infty$ with $K$ fixed.

Our goal is not to claim a complete theory at this stage, but rather to isolate the objects that
should play the role of the WABC limits in our hierarchical setting, and to explain why the
large-$M$ over-sampling posterior behaves very differently from the target \textsc{permABC}
pseudo-posterior.

\subsection{A compartment-level Wasserstein viewpoint}

Recall that the observed data are
\[
  y = (y_1,\dots,y_K) \in \mathcal D^K,
\]
where each compartment $y_k \in \mathcal D$ depends on a global parameter $\beta$ and a local
parameter $\mu_k$. For simulated data
\[
  z = (z_1,\dots,z_K) \in \mathcal D^K,
\]
our acceptance criterion is based on
\[
  \min_{\sigma \in S_K} d(y,z_\sigma).
\]

It is useful to introduce the empirical measures on the space of compartments,
\[
  \widehat\nu_y^K := \frac1K \sum_{k=1}^K \delta_{y_k},
  \qquad
  \widehat\nu_z^K := \frac1K \sum_{k=1}^K \delta_{z_k}.
\]
If the compartment-level discrepancy is induced by a ground cost $\rho$ on $\mathcal D$, i.e.
\[
  d_K(y,z)^p := \frac1K \sum_{k=1}^K \rho(y_k,z_k)^p,
\]
then
\[
  \min_{\sigma\in S_K} d_K(y,z_\sigma)
  = W_p(\widehat\nu_y^K,\widehat\nu_z^K),
\]
where $W_p$ is the Wasserstein distance on $\mathcal P_p(\mathcal D)$. If one uses instead an
unnormalized distance, then the relation still holds up to the deterministic factor $K^{1/p}$.

Thus, once the data are viewed at the \emph{compartment level}, \textsc{permABC} is the natural
hierarchical analogue of WABC. The difference is that WABC works directly on an exchangeable
sample $y_{1:K}$, whereas \textsc{permABC} also carries a latent representation
\[
  \theta = (\beta,\mu_1,\dots,\mu_K),
\]
for which the compartment labels matter. The role of the projection step is precisely to recover
an ordered local configuration after solving the transport/matching problem.

\begin{remark}[Dictionary with WABC]
Ignoring the local labels, the permutation-invariant pseudo-posterior $\tilde\pi_\varepsilon$
behaves like WABC applied to the empirical distribution of compartments. By contrast, the
projected pseudo-posterior $\pi_\varepsilon^\star$ contains extra information: it keeps track of
which local parameters are attached to the optimally matched compartments.
\end{remark}

\subsection{A WABC-style limit when $K \to \infty$}

The relevant object when $K$ grows is not the full joint law of the labeled vector
$(y_1,\dots,y_K)$, but the \emph{marginal law of one compartment}. Define
\[
  \nu_\beta(\mathrm dy)
  := \int g(\mathrm dy\mid \beta,\mu)\,\pi_\mu(\mathrm d\mu\mid \beta),
\]
where the conditional prior notation $\pi_\mu(\cdot\mid \beta)$ is only used for generality; in
many examples it reduces to a fixed prior $\pi_\mu$.

Under the model, conditional on $\beta$, the compartments are exchangeable, and their empirical
measure should converge to $\nu_\beta$. Likewise, if the observed data come from some true
mechanism, we write $\nu_\star$ for the large-$K$ limit of $\widehat\nu_y^K$.

This suggests that the correct large-$K$ asymptotic target for the \emph{global} parameter is
\[
  \beta^\star \in \arg\min_{\beta} W_p(\nu_\beta,\nu_\star),
\]
exactly in the spirit of WABC.

\begin{proposition}[Heuristic WABC analogue for $K\to\infty$]
Assume that:
\begin{enumerate}
  \item $\widehat\nu_y^K \Rightarrow \nu_\star$ as $K\to\infty$;
  \item for each fixed $\beta$, the simulated empirical measure satisfies
  $\widehat\nu_z^K \Rightarrow \nu_\beta$ in probability under the prior predictive model;
  \item the map $\beta \mapsto \nu_\beta$ is identifiable in Wasserstein distance.
\end{enumerate}
Then the permutation-invariant pseudo-posterior should behave as a coarsened posterior on
$\beta$, namely
\[
  \tilde\pi_{\varepsilon_K}(\mathrm d\beta\mid y)
  \approx
  \pi_\beta(\mathrm d\beta)\,
  \mathbf 1\!\left\{W_p(\nu_\beta,\nu_\star) \lesssim \varepsilon_K\right\},
\]
up to the empirical fluctuations of $\widehat\nu_y^K$ and $\widehat\nu_z^K$.
In particular, if $\varepsilon_K \to 0$ slowly enough, one expects concentration near the set of
minimizers of $\beta \mapsto W_p(\nu_\beta,\nu_\star)$.
\end{proposition}

\begin{proof}[Informal justification]
The acceptance event only depends on the optimally matched discrepancy between the two
empirical measures of compartments. When $K$ is large, both empirical measures should be close
to their population limits. Therefore the acceptance event becomes asymptotically equivalent to
\[
  W_p(\nu_\beta,\nu_\star) \lesssim \varepsilon_K,
\]
which is exactly the WABC mechanism after replacing the original parameter $\theta$ by the
integrated compartment law $\nu_\beta$.
\end{proof}

\begin{remark}[What this says --- and what it does not say]
This asymptotic picture is naturally a statement about $\beta$, not about the labeled vector
$\mu_{1:K}$. The reason is that transport over compartments forgets labels. Therefore the WABC
analogy is fundamentally a statement about the integrated model $\beta \mapsto \nu_\beta$.
\end{remark}

\subsection{Why the local pseudo-posterior is not a usual Bayesian posterior}

The previous discussion suggests a clean WABC-like limit for the global parameter, but not for
the local parameters. This is not an accident.

Indeed, after permutation, the local parameters that appear in $\pi_\varepsilon^\star$ are not the
original latent variables of the generative model in a labeled sense; rather, they are the local
parameters \emph{selected by the optimal matching}. Hence they should be interpreted as
\emph{transport witnesses} or \emph{best local explanations} of the observed compartments, not as
a standard posterior on pre-existing labeled latent variables.

For this reason, even when $K \to \infty$, the asymptotic behaviour of the local part of
$\pi_\varepsilon^\star$ should be expected to differ from ordinary posterior consistency. At best, one
might hope for convergence of the empirical law of matched locals, or for convergence toward a
transport coupling between observed compartments and simulated latent explanations.

\subsection{Large-$M$ over-sampling is a different asymptotic object}

We now turn to over-sampling. Recall that for $M > K$ we simulate
\[
  \theta = (\beta,\mu_1,\dots,\mu_M),
  \qquad
  z=(z_1,\dots,z_M),
\]
and accept if
\[
  \min_{\sigma\in \mathcal A_K^M} d(y,z_\sigma) \le \varepsilon,
\]
where $\mathcal A_K^M$ denotes the set of injections from $\{1,\dots,K\}$ into
$\{1,\dots,M\}$.

This defines two different objects:
\begin{itemize}
  \item the \emph{unprojected} over-sampling pseudo-posterior $\tilde\pi^{M}_{\varepsilon}$ on
  $(\beta,\mu_{1:M})$;
  \item the \emph{projected} pseudo-posterior $\pi^{\star,M}_{\varepsilon}$ obtained after keeping the
  $K$ matched components selected by the optimal injection.
\end{itemize}
These two objects have very different large-$M$ limits.

\begin{proposition}[Heuristic large-$M$ behaviour of the unprojected law]
Fix $K$ and $\varepsilon>0$. Suppose that for a given $\beta$ and for every observed compartment
$y_k$, the probability that one simulated compartment falls in an $\varepsilon$-matching region of
$y_k$ is positive. Then the acceptance probability under over-sampling tends to one as
$M\to\infty$. Consequently, the $\beta$-marginal of the unprojected pseudo-posterior satisfies
\[
  \tilde\pi^{M}_{\varepsilon}(\mathrm d\beta\mid y)
  \longrightarrow
  \pi_\beta(\mathrm d\beta),
\]
or more generally to the prior restricted to the set of $\beta$ for which matching is feasible.
\end{proposition}

\begin{proof}[Informal justification]
For each $y_k$, let $B_{k,\varepsilon}(\beta)$ denote the event that one simulated compartment is
close enough to $y_k$ to be part of a valid matching. Under independence across the $M$
simulated compartments, if each event has positive probability, then the probability that none of
the $M$ candidates matches $y_k$ decays exponentially in $M$. Hence the probability that one can
find $K$ distinct matches tends to one. Once acceptance becomes almost certain, the ABC filter
no longer discriminates between values of $\beta$, so the marginal over $\beta$ reverts to its prior.
\end{proof}

This explains the phenomenon already observed numerically: when $M$ is very large, the
unprojected over-sampling pseudo-posterior is not an approximation of the target
$\textsc{permABC}$ pseudo-posterior at all. It is instead an \emph{auxiliary bridge} which becomes
progressively less informative about $\beta$ as $M$ grows.

\subsection{Why the projected local law becomes concentrated when $M \to \infty$}

The projected law $\pi^{\star,M}_{\varepsilon}$ behaves very differently. Although the global marginal
becomes almost prior-like, the selected local components are chosen among $M$ candidates by an
extreme-value mechanism: they are the \emph{best matches} among an increasingly large pool.

To formalize this, fix $(\beta,y_k)$ and let
\[
  c_\beta(y_k,\mu)
\]
denote a local mismatch criterion. Depending on the context, one may take for example
\[
  c_\beta(y_k,\mu) = d\bigl(y_k,G(\beta,\mu)\bigr)
\]
in a deterministic local simulator $z=G(\beta,\mu)$, or an expected/marginal discrepancy in the
stochastic case.

If the $\mu_j$ are i.i.d. from $\pi_\mu(\cdot\mid \beta)$, then the matched locals in the projected
law are approximately the minimizers of $c_\beta(y_k,\mu_j)$ over $j=1,\dots,M$. Hence they are
order-statistic-type variables.

\begin{proposition}[Heuristic concentration of projected locals as $M\to\infty$]
Assume that for each $(\beta,y_k)$ the criterion $\mu \mapsto c_\beta(y_k,\mu)$ has a unique
minimizer
\[
  \mu_k^{\dagger}(\beta,y_k)
  \in \arg\min_\mu c_\beta(y_k,\mu),
\]
and that this minimizer lies in the support of the local prior. Then under the projected
large-$M$ over-sampling law one expects
\[
  \mu_k^{\star,M} \Longrightarrow \delta_{\mu_k^{\dagger}(\beta,y_k)}
\qquad \text{conditionally on } \beta,
\]
for each matched component $k=1,\dots,K$.
\end{proposition}

In words: the selected local pseudo-posterior becomes increasingly concentrated around the best
possible local explanation of each observed compartment under the current global parameter.
This is why the large-$M$ local law can look ``Dirac-like'' even though the global marginal in
$\beta$ stays essentially equal to the prior.

\begin{remark}[Deterministic and invertible local model]
In the particularly transparent case where $z_k = G(\beta,\mu_k)$ deterministically and the map
$\mu \mapsto G(\beta,\mu)$ is injective, the best local explanation is simply the inverse image of
$y_k$, whenever it exists. In that case the large-$M$ projected local law really does converge to a
Dirac mass at the inverse solution.
\end{remark}

\subsection{A non-commutation of limits}

The previous discussion highlights a key point:
\[
  \lim_{M\to\infty} \tilde\pi^{M}_{\varepsilon}(\mathrm d\beta\mid y)
  = \pi_\beta(\mathrm d\beta)
\]
in general, whereas the target distribution is obtained at
\[
  M = K,
  \qquad
  \tilde\pi^{K}_{\varepsilon} = \tilde\pi_{\varepsilon}.
\]
Therefore the limit $M\to\infty$ is \emph{not} the inferential target. It is only useful because it
creates an easier auxiliary distribution from which an SMC algorithm can start. This also explains
why the schedule $M_0 > M_1 > \cdots > M_T = K$ is essential: the statistical meaning is only
recovered at the end of the bridge.

\subsection{Take-home messages}

The discussion above suggests the following interpretation.
\begin{enumerate}
  \item At the level of \emph{empirical distributions of compartments}, \textsc{permABC} is the
  natural hierarchical analogue of WABC.
  \item When $K\to\infty$, the relevant asymptotic object for the global parameter is the
  integrated compartment law $\nu_\beta$, and one expects WABC-style concentration toward
  minimizers of $\beta \mapsto W_p(\nu_\beta,\nu_\star)$.
  \item This asymptotic picture does \emph{not} imply ordinary posterior consistency for the
  labeled locals $\mu_{1:K}$, since the projection step turns them into matched witnesses rather
  than genuine labeled latent variables.
  \item When $M\to\infty$ with $K$ fixed, over-sampling produces a qualitatively different
  object: the global marginal becomes prior-like, while the projected locals become extreme-value
  or best-match objects, often close to Dirac masses.
  \item Consequently, over-sampling should be viewed as an algorithmic bridge, not as an
  inferential target in its own right.
\end{enumerate}

\subsection{Open questions}

Several natural questions remain open.
\begin{itemize}
  \item Can one prove a full WABC-style concentration result for $\beta$ under $K\to\infty$,
  with explicit rates?
  \item What is the correct asymptotic object for the empirical distribution of matched local
  parameters?
  \item What happens in a joint regime $K\to\infty$ and $M\to\infty$, for instance when
  $M/K \to c>1$?
  \item Can the projected local limit be characterized as a transport map, a Monge coupling, or
  a nearest-neighbour selection rule?
  \item Which assumptions guarantee that the large-$M$ projected local law converges to a Dirac
  mass, and when does it instead converge to a non-degenerate set of minimizers?
\end{itemize}

We view these questions as promising directions for extending the current theoretical picture of
\textsc{permABC} beyond the fixed-$K$ guarantees already established in this paper.
